\documentclass[10pt,a4paper]{article}
\usepackage[utf8]{inputenc}
\usepackage[vietnamese]{babel}
\usepackage{amsmath,amssymb}
\usepackage{mathptmx}
\usepackage[top=20mm, bottom=20mm, left=18mm, right=18mm, headheight=14pt]{geometry}
\usepackage{graphicx}
\usepackage{xcolor}
\usepackage{tikz}
\usepackage{fancyhdr}

% Định nghĩa màu sắc nhận diện thương hiệu Stewart Calculus
\definecolor{stewartcyan}{RGB}{0, 130, 185}
\definecolor{stewartred}{RGB}{204, 30, 30}
\definecolor{stewartpurple}{RGB}{112, 48, 160}
\definecolor{stewartlightpurple}{RGB}{237, 233, 246}

% Thiết lập tiêu đề đầu trang
\pagestyle{fancy}
\fancyhf{}
\fancyhead[R]{\fontsize{8.5}{11}\selectfont\textsf{\textbf{MỤC 3.3}}\hspace{1.5em}\textrm{Đạo hàm của các hàm lượng giác}\hspace{2.5em}\textsf{\textbf{195}}}
\renewcommand{\headrulewidth}{0pt}

\setlength{\parindent}{14pt}
\setlength{\parskip}{0pt}

\begin{document}

\noindent
\begin{minipage}[t]{0.26\textwidth}
  \vspace{0pt}
  \centering
  \includegraphics[width=\linewidth]{images/p0230_fig1.png}
  
  \vspace{3pt}
  {\small\textsf{\textbf{HÌNH 5}}}
  
  \vspace{2.5cm}
  \noindent
  \begin{tikzpicture}
    \node[fill=stewartlightpurple, inner sep=0pt, anchor=west] {
      \colorbox{stewartpurple}{\rule[-2pt]{0pt}{9.5pt}\hspace{2.5pt}\textbf{\textsf{\footnotesize\color{white}PS}}\hspace{2.5pt}}%
      \hspace{4pt}\textbf{\textsf{\scriptsize Tìm kiếm quy luật.}}\hspace{6pt}
    };
  \end{tikzpicture}
  
  \vspace{4.7cm}
  \includegraphics[width=\linewidth]{images/p0230_fig2.png}
  
  \vspace{3pt}
  {\small\textsf{\textbf{HÌNH 6}}}
\end{minipage}%
\hfill
\begin{minipage}[t]{0.71\textwidth}
  \vspace{0pt}
  \noindent Vật dao động từ vị trí thấp nhất ($s = -4\text{ cm}$) đến vị trí cao nhất ($s = 4\text{ cm}$). Chu kỳ của dao động là $2\pi$, bằng chu kỳ của $\cos t$.
  
  \indent Tốc độ là $|v| = 4|\sin t|$, đạt giá trị lớn nhất khi $|\sin t| = 1$, tức là khi $\cos t = 0$. Do đó, vật chuyển động nhanh nhất khi nó đi qua vị trí cân bằng ($s = 0$). Tốc độ của nó bằng $0$ khi $\sin t = 0$, nghĩa là tại các vị trí cao nhất và thấp nhất.
  
  \indent Gia tốc $a = -4\cos t = 0$ khi $s = 0$. Nó có độ lớn cực đại tại các điểm cao nhất và thấp nhất. Xem các đồ thị ở Hình~5.\hfill{\color{stewartcyan}$\blacksquare$}

  \vspace{1.1em}
  \noindent
  \textbf{\textsf{\color{stewartcyan}VÍ DỤ 4}} \quad Tìm đạo hàm cấp 27 của $\cos x$.

  \vspace{0.4em}
  \noindent
  \textbf{\textsf{\color{stewartcyan}LỜI GIẢI}} \quad Một vài đạo hàm đầu tiên của hàm số $f(x) = \cos x$ như sau:
  \begin{align*}
    f'(x) &= -\sin x \\
    f''(x) &= -\cos x \\
    f'''(x) &= \sin x \\
    f^{(4)}(x) &= \cos x \\
    f^{(5)}(x) &= -\sin x
  \end{align*}
  Ta thấy rằng các đạo hàm liên tiếp lặp lại theo một chu kỳ có độ dài bằng 4 và đặc biệt là, $f^{(n)}(x) = \cos x$ mỗi khi $n$ là một bội số của 4. Do đó
  \[
    f^{(24)}(x) = \cos x
  \]
  và, lấy đạo hàm thêm ba lần nữa, ta có
  \begin{equation*}
    f^{(27)}(x) = \sin x \tag*{{\color{stewartcyan}$\blacksquare$}}
  \end{equation*}

  \vspace{0.8em}
  \noindent
  {\color{stewartcyan}\rule{1.1ex}{1.1ex}}\hspace{0.5em}\textbf{\textsf{Hai giới hạn lượng giác đặc biệt}}

  \vspace{0.5em}
  \noindent
  Khi chứng minh công thức tính đạo hàm của hàm sin, chúng ta đã sử dụng hai giới hạn đặc biệt mà bây giờ chúng ta sẽ chứng minh.

  \vspace{0.5em}
  \noindent
  \begin{minipage}{\linewidth}
    \makebox[0pt][l]{%
      \tikz[baseline=(box5.base)]{%
        \node[draw=stewartred, line width=1pt, rounded corners=0.5pt, inner sep=3pt] (box5) {\bfseries\sffamily\color{stewartred}\small 5};
      }%
    }%
    \centering
    \tikz[baseline=(eq.base)]{%
      \node[draw=stewartred, line width=1.1pt, inner xsep=22pt, inner ysep=7pt] (eq) {$\displaystyle \lim_{\theta \to 0} \frac{\sin \theta}{\theta} = 1$};
    }%
  \end{minipage}
  \vspace{0.6em}

  \noindent
  \textbf{\textsf{\color{stewartred}CHỨNG MINH}} \quad Trước hết, giả sử rằng $\theta$ nằm giữa $0$ và $\pi/2$. Hình~6(a) biểu diễn một hình quạt tròn tâm $O$, góc ở tâm $\theta$, và bán kính bằng 1. Đoạn thẳng $BC$ được vẽ vuông góc với $OA$. Theo định nghĩa của số đo radian, ta có $\text{cung } AB = \theta$. Ta cũng có $|BC| = |OB|\sin\theta = \sin\theta$. Từ hình vẽ, ta thấy rằng
  \[
    |BC| < |AB| < \text{cung } AB
  \]
  \noindent
  Do đó \hfill $\displaystyle \sin\theta < \theta \qquad \text{vì vậy} \qquad \frac{\sin\theta}{\theta} < 1$ \hfill\null

  \vspace{0.5em}
  \noindent
  Gọi giao điểm của các tiếp tuyến tại $A$ và $B$ là $E$. Bạn có thể thấy từ Hình~6(b) rằng chu vi của một đường tròn nhỏ hơn chu vi của đa giác ngoại tiếp, và do đó $\text{cung } AB < |AE| + |EB|$. Như vậy
  \begin{align*}
    \theta = \text{cung } AB &< |AE| + |EB| \\
    &< |AE| + |ED| = |AD| = |OA|\tan\theta \\
    &= \tan\theta
  \end{align*}
\end{minipage}

\end{document}